\chapter{Конструкторский раздел}

В этом разделе будут рассмотрены алгоритмы, выбранные для построения изображения, а также спроектированна диаграмма классов приложения.

\section{Алгоритм построения изображений}
В качестве средства описания алгоритма построения изображения будет использоваться диаграмма IDEF0.

Построение изображений в компьютерной графике в контексте диаграммы IDEF0 определяется следующими параметрами:

\noindent \textbf{Описание входных данных}
\begin{itemize}
    \item Объекты определяются геометрией, положением и ориентацией и свойствами материала (цвет, текстуры, нормали);
    \item Камера определяется проекцинными параметрами (перспективная, ортографическая), положением и ориентацией;
    \item Источники освещения определяются типом (направленный, точечный и др.) , положением и ориентацией (для направленных), цветом и интенсивностью.
\end{itemize}

\noindent \textbf{Описание элементов управления}
\begin{itemize}
    \item Алгоритм рендеринга: растеризация, трассировка лучей, трассировка пути, ray marching;
    \item Параметры рендеринга: разрешение получаемого изображения, максимальное число переотражений луча, количество сэмплов на пиксель;
    \item Модель освещения: локальная (Фонг, Блинн-Фонг) или глобальная (PBR\cite{pbr_book});
    \item Способ затенения: простое, по Гуро, по Фонгу.
\end{itemize}

\noindent \textbf{Описание механизмов}
\begin{itemize}
    \item Аппаратное обеспечение (CPU, RAM, GPU + VRAM);
    \item Графический API (OpenGL, Vulkan, DirectX, Metal и др.);
    \item Библиотека для окна и ввода (GLFW, SDL, SFML и др.).
    \item Язык программирования (C++, Python, C\#, Java и др.);
\end{itemize}

\noindent \textbf{Выходными данными} является сгенерированный кадр.

Общий алгоритм генерация изображений представлен на рисунке \ref{fig:IDEF0_general}.
\begin{figure}[H]
    \centering
    \includegraphics[width=1\textwidth]{IDEF0_general.pdf}
    \caption{IDEF0 диаграмма общего алгоритма построения изображения}
    \label{fig:IDEF0_general}
\end{figure}

Выбранный для реализации алгоритм построения изображений с конкретными методами, моделями и инструментами представлен на рисунке \ref{fig:IDEF0_specific}.
\begin{figure}[H]
    \centering
    \includegraphics[width=1\textwidth]{IDEF0_specific.pdf}
    \caption{IDEF0 диаграмма верхнего уровня (A-0) выбранного для реализации алгоритма построения изображения}
    \label{fig:IDEF0_specific}
\end{figure}


\section{Алгоритм обратной трассировки лучей}
IDEF0 диаграмма уровня A-1 алгоритма трассировки лучей представлена на рисунке \ref{fig:IDEF0_trace_ray}.
\begin{figure}[H]
    \centering
    \includegraphics[width=0.8\textwidth]{IDEF0_RayTracing.pdf}
    \caption{IDEF0 диаграмма трассировки лучей}
    \label{fig:IDEF0_trace_ray}
\end{figure}

Схема алгоритма обратной трассироки лучей приведена на рисунке \ref{fig:trace_ray_diagram}.
\begin{figure}[H]
    \centering
    \includegraphics[width=0.6\textwidth]{trace_ray_diagram.pdf}
    \caption{Схема алгоритма обратной трассировки лучей}
    \label{fig:trace_ray_diagram}
\end{figure}